\practiceweek{1}

\begin{exercisebox}
\begin{pproblem}[Membership needs $n$ queries]\label{pp:membership}
    The scribe shows that finding $x$ needs $n - 1$ queries when $x$ is promised to be in the array.
    Show that deciding whether $x$ is in the array at all needs $n$ queries.
\end{pproblem}
\tcblower
\begin{wsolution}
    \uses{\scr{thm:detsearch-0825}}
    \leavevmode
    \begin{enumerate}
        \item Fix a deterministic algorithm. The adversary answers every query with ``not $x$''.
        \item Suppose the algorithm stops after at most $n - 1$ queries. Then some slot $j$ was never
              queried.
        \item Two arrays agree with all the answers: one has $x$ in slot $j$, the other has no $x$
              anywhere.
        \item The algorithm saw the same answers on both, so it gives the same output on both. It is
              wrong on one.
        \item So a correct algorithm makes $n$ queries.
    \end{enumerate}
    The promise is what saved one query on Aug 25: with a promise, $n - 1$ answers of ``no'' leave
    only one possible slot.
\end{wsolution}
\end{exercisebox}

\begin{exercisebox}
\begin{pproblem}[Search with $k$ copies of $x$]\label{wp:kcopies}
    An array of $n$ slots contains $x$ in exactly $k$ slots, $1 \leq k < n$.
    \begin{enumerate}[label=(\alph*)]
        \item Random sampling queries uniformly random slots (with replacement) until it finds $x$.
              What is the expected number of queries?
        \item Show that a deterministic algorithm needs $n - k$ queries in the worst case.
    \end{enumerate}
\end{pproblem}
\tcblower
\begin{wsolution}
    \uses{\scr{thm:geometric-0827}; \scr{thm:detsearch-0825}}
    (a) Each query finds $x$ with probability $p = \sfrac{k}{n}$, independently of the others, so the
    number of queries is geometric with $p = \sfrac{k}{n}$. Its expectation is
    $\sfrac{1}{p} = \sfrac{n}{k}$.

    (b) The adversary answers ``not $x$'' to every query. Suppose the algorithm stops after
    $q < n - k$ queries and outputs a slot $j$. The array has $n - k$ slots without $x$, and only
    $q$ of them have been used up by the answers, so at least one more slot without $x$ remains. If
    $j$ was queried, the answer at $j$ was ``not $x$'' and the output is wrong. Otherwise the
    adversary puts the remaining empty slot at $j$ and the $k$ copies of $x$ in other unqueried
    slots; this agrees with all the answers, and the output is wrong. After $n - k$ answers of
    ``no'', the $k$ unqueried slots all contain $x$, so $n - k$ queries suffice.

    For $k = 1$ this is the $n - 1$ of the scribe.
\end{wsolution}
\end{exercisebox}
